%META {"title":"Cùng vuông góc với đường thẳng thứ ba","note":"a ⊥ c và b ⊥ c (hai dấu góc vuông) → a // b. HS trả lời miệng: hai đường thẳng phân biệt cùng vuông góc với một đường thẳng thứ ba thì song song với nhau."}
\documentclass[border=5pt]{standalone}
\usepackage{tkz-euclide}
% Màu dùng chung cho toàn bộ hình vẽ (ảnh stack với nền tối của web)
\definecolor{tminc}{RGB}{226,232,240}   % chữ + nét chính
\definecolor{tmblue}{RGB}{0,210,255}    % nhấn neon xanh
\definecolor{tmpink}{RGB}{255,0,200}    % nhấn neon hồng
\definecolor{tmgreen}{RGB}{0,255,136}   % nhấn neon xanh lá
\color{tminc}

\begin{document}
\begin{tikzpicture}[line width=1.1pt]
  \coordinate (A) at (0.9,-1);
  \coordinate (B) at (0.9,2.5);
  \coordinate (C) at (3.1,-1);
  \coordinate (D) at (3.1,2.5);
  \coordinate (E) at (-0.5,0);
  \coordinate (F) at (4,0);
  \coordinate (I) at (0.9,0);
  \coordinate (J) at (3.1,0);
  \coordinate (U) at (0.9,0.6);
  \coordinate (R) at (1.5,0);
  \coordinate (V) at (3.1,0.6);
  \coordinate (S) at (2.5,0);
  \draw[tmblue] (A) -- (B) node[above=2pt]{$a$};
  \draw[tmblue] (C) -- (D) node[above=2pt]{$b$};
  \draw[tminc] (E) -- (F) node[right=2pt]{$c$};
  \tkzMarkRightAngles[size=0.3,color=tmpink,line width=1pt](U,I,R V,J,S)
\end{tikzpicture}
\end{document}
